In this paper, we developed a formal theory of convergence detection to address the fundamental gap between the mathematical specifications and practical implementations of Fixpoint Detection (FPD) in modern incremental recursive computation frameworks like DBSP. We identified and formalized the FPD problem, demonstrated that the commonly suggested ``FirstZero'' heuristic is unsound for the ``delta-of-deltas'' incremental optimization for recursive computation, and proved that exact FPD is undecidable for arbitrary expressive circuits. Within our theory, we defined internal convergence (IntConv) as a declarative criterion corresponding to the internal-state-stability strategy used by practical implementations, and proved that it is a sound sufficient condition for convergence. We also defined the State Fixpoint (StFP) predicate and a sound and complete StFP detector, providing semantic and algorithmic formalizations of this strategy. The sound and complete global IntConv detector combines the StFP detector with fixed-input and zero-output checks across all bracketed sub-circuits. These predicates and detectors treat nested bracketed circuits compositionally. Furthermore, we proved that IntConv is complete for a broad class of useful circuits, including arbitrarily nested while loops, Datalog queries, and their incrementally optimized versions. These results provide formal semantic guarantees for an existing implementation strategy, and our results have been formalized in Lean 4.

Our results are mainly theoretical, rather than an empirical performance evaluation. A natural direction for future work is to implement and compare the global IntConv detector based on the general StFP detector with specialized global detectors within a complete system. Such an evaluation should measure detector overhead and compare end-to-end incremental execution with non-incremental recomputation, since that comparison also depends on the query, workload, data representation, and other system-level optimizations.
